\clearpage
\section{References}

\begin{multicols}{2}
\footnotesize
\RaggedRight

\subsection{Scripture and doctrine}

\begin{itemize}
\item Sacred Scripture: Ex 3; 19; 20; 24; 1 Kgs 17--19; 2 Kgs 2; Ps 24; 62; 1 Chr 28:9; Sir 48; Is 35; 40:1--11; 42; 61; Jer 2; 31; Hos 2; Mic 6:6--8; Mal 3--4; Mt 4--6; 11; 17; Lk 1--2; 4; 9:18--25, 51--56; 10:38--42; Jn 2:1--12; 10:7--9; 19:25--27; Acts 1:14; 2:42--47; 4:32--35; Rom 6; 13:8--14; 1 Cor 12; Gal 3:23--29; Col 3:1--17; Heb 12:18--29; Jas 3; Rev 2:10; 21.
\item Second Vatican Council, \textit{Lumen gentium} 53, 56--62, 66--69; \textit{Sacrosanctum Concilium} 13; \textit{Dei Verbum} 8--10, 21--26.
\item Dicastery for the Doctrine of the Faith, \textit{Mater Populi Fidelis}: Doctrinal Note on Some Marian Titles Regarding Mary's Cooperation in the Work of Salvation (4 November 2025), especially nos.~27--33, 45--55, and 62--70, \sourceurl{https://www.vatican.va/roman_curia/congregations/cfaith/documents/rc_ddf_doc_20251104_mater-populi-fidelis_en.html}{official Holy See text}; used for the current limits on mediation, intercession, and the communication of grace.
\item \textit{Catechism of the Catholic Church}, nos.~487--507, 963--975, 1667--1673, 2016, 2673--2679, on Mary, sacramentals, perseverance, and prayer.
\item St.~Thomas Aquinas, \textit{Summa theologiae} III, qq.~27--30; III, q.~60, aa.~1--5; Supplement, q.~71, on Marian grace, sacred signs, and suffrage, used within the limits stated in the text.
\end{itemize}

\subsection{Carmelite origins}

\begin{itemize}
\item Rule of Saint Albert, received between 1206 and 1214, especially nos.~2, 7, 10--12, 14--24, \sourceurl{https://carmelite.org/spirituality/rule-of-saint-albert/}{Carmelite text and historical introduction}.
\item Honorius III, \textit{Ut vivendi normam} (30 January 1226); Gregory IX confirmation (1229); Innocent IV, \textit{Quae honorem Conditoris} (1 October 1247). The Carmelite witnesses distinguish the original \textit{formula vitae} from the adapted bullated Rule; see the \sourceurl{https://www.ordem-do-carmo.pt/index.php/portal/1403-8-centenario-da-aprovacao-da-regra-original-de-santo-alberto-e-da-ordem-carmelita-pelo-papa-honorio-iii-30-de-janeiro-de-2026}{Carmelite eighth-centenary source record}.
\item Carmelite Constitutions and formation texts on allegiance to Jesus Christ, prayer, fraternity, Mary, Elijah, service, and discretion.
\item Joachim Smet, \textit{The Carmelites: A History of the Brothers of Our Lady of Mount Carmel}, used as historical synthesis rather than magisterial authority.
\end{itemize}

\subsection{Feast and Scapular}

\begin{itemize}
\item Congregation for Divine Worship, \textit{Collectio Missarum de Beata Maria Virgine}, editio typica (Libreria Editrice Vaticana, 1987), decree Prot. N. 309/86 (15 August 1986). John Paul II approved the collection and ordered its publication; the decree describes it as an appendix to the Roman Missal, and Praenotanda 22 says its promulgation introduces no change to the 1975 \textit{Missale Romanum} or the governing rubrics. The nine daily collects are from formularies 10 (p.~42), 20 (p.~83), 22 (p.~90), 23 (p.~93), 25 (p.~100), 28 (p.~112), 38 (p.~148), 41 (p.~159), and 46 (p.~175). Formulary 32 and its collect on pp.~126--127 supply the Christ-as-mountain theme but are not assigned as a recited prayer. See the Dicastery's \sourceurl{https://www.cultodivino.va/en/formazione/pubblicazioni/libri-liturgici/aliae/collectio-missarum-de-beata-maria-virgine.html}{official edition record} and the \sourceurl{https://www.liturgia.it/content/BMV/Collectio\%20Missarum\%20de\%20Beata\%20Maria\%20Virgine\%20\%281987\%29.pdf}{complete Latin facsimile}.
\item \textit{General Instruction of the Roman Missal}, no.~54, \sourceurl{https://www.vatican.va/content/dam/wss/roman_curia/congregations/ccdds/documents/rc_con_ccdds_doc_20030317_ordinamento-messale_en.html}{official Holy See text}. It supplies the exact longer Roman-Missal endings represented by \textit{Per Dominum.} and \textit{Qui tecum vivit.}, together with the people's \textit{Amen}. Those governing formulas, not an editorial composition, complete the nine abbreviated conclusions.
\item \textit{Collection of Masses of the Blessed Virgin Mary}, vol.~I, \textit{Sacramentary} (Catholic Book Publishing, 2012), formularies 10, 20, 22, 23, 25, 28, 38, 41, and 46. The current approved English was checked but is cited rather than reproduced. ICEL's standing \sourceurl{https://www.icelweb.org/copyright.htm}{copyright policy} requires written permission or a license, which was not obtained. The USCCB's published successor \sourceurl{https://www.usccb.org/committees/divine-worship/policies/guidelines-for-the-publication-of-liturgical-books/devotional-and-other-non-liturgical-publications}{devotional-publication guidance}, approved in 2025 and effective 29 November 2026, likewise requires authentication and permission; it is identified as a forthcoming rule, not one already in force on this guide's as-of date.
\item \textit{Missale Romanum} (1962), Good Friday order no.~960, printed p.~180, for the shared Lord's Prayer Latin; and \textit{Ordo Missae}, printed p.~216, for the shared Lesser Doxology Latin. Rubrical recitation marks are omitted from the Lord's Prayer, but its words, punctuation, stress accents, ligatures, and final response are retained.
\item Dom Gaspar Lefebvre, O.S.B., \textit{Daily Missal with Vespers} (St.~Andrew's Abbey; E.~M. Lohmann, 1925), pp.~4--5 for the historical English of the shared Lord's Prayer and Lesser Doxology; p.~35 for the Sign of the Cross in Latin and English; p.~5 for the Hail Mary in Latin and English; pp.~1499--1500 for the recited 16 July collect in Latin and English; and pp.~31--32 for its full \textit{Qui vivis} conclusion. The volume bears the 13 May 1925 imprimatur of Archbishop Austin Dowling and is public domain in the United States.
\item Congregation for Divine Worship and the Discipline of the Sacraments, \textit{Directory on Popular Piety and the Liturgy} (2001), nos.~183--189, 204--205, \sourceurl{https://www.vatican.va/roman_curia/congregations/ccdds/documents/rc_con_ccdds_doc_20020513_vers-direttorio_en.html}{official English text}.
\item \textit{Rite of Blessing of and Enrollment in the Scapular of the Blessed Virgin Mary of Mount Carmel}, approved by the same dicastery on 5 January 1996. Its formula and theology are reported, not reproduced as this booklet's rite.
\item Pius XII, \textit{Neminem profecto latet} (11 February 1950), AAS 42 (1950), 390--391.
\item John Paul II, Angelus address (24 July 1988), and Message to the Carmelite Family (25 March 2001), \sourceurl{https://www.vatican.va/content/john-paul-ii/en/messages/pont_messages/2001/documents/hf_jp-ii_mes_20010326_ordine-carmelo.html}{official English text}.
\item Benedict XIII's extension of the feast to the Latin Church on 24 September 1726; current General Roman Calendar and Carmelite proper-calendar status, checked through 13 July 2026; see the \sourceurl{https://www.carmelites.org.au/item/1722-feast-day-message-from-prior-general}{Carmelite Prior General's 2026 historical notice}.
\end{itemize}

\subsection{Historical criticism}

\begin{itemize}
\item Richard Copsey, O.Carm., ``Simon Stock and the Scapular Vision,'' \textit{Journal of Ecclesiastical History} 50 (1999): 652--683, \sourceurl{https://doi.org/10.1017/S002204699900250X}{article record}; used for the late and disputed documentary provenance.
\item Sam Anthony Morello, OCD, and Patrick McMahon, O.Carm., ``Catechesis on the Brown Scapular,'' \sourceurl{https://www.carmelitefriarsocd.org/blog/brown-scapular}{Carmelite provincial text}; used for the sacramental, affiliation, Simon, and Sabbatine boundaries.
\item Carmelite historical account of the Marian commemoration's emergence and 16 July date, \sourceurl{https://www.carmelitepriory.org/origins-solemn-commemoration-lady-of-mount-carmel/}{Carmelite Priory study}; conclusions retained at their stated historical level.
\item \textit{Flos Carmeli} reception and hymnological witnesses, including John Sullivan, ``Liturgical Creativity from Edith Stein,'' \textit{Teresianum} 49 (1998): 168--170 and nn.~12, 15--16, \sourceurl{https://www.teresianum.net/wp-content/uploads/2016/11/Ter_49_1998-1_165-185.pdf}{article}; and the \sourceurl{https://ocdssacramento.org/wp-content/uploads/2019/07/office2007.pdf}{2007 \textit{Discalced Carmelite Proper Offices}}, pp.~61--63, for the complete singular \textit{hic adesse / me tibi servulum} witness and the familiar complete-English alternative. The sequence has a medieval core and later expansion. Simon Stock authorship, textual growth, and the exact earliest form remain uncertain; \textit{esto propitia} and \textit{da privilegia}, and \textit{hic adesse} and \textit{nos ad esse}, are recorded variants.
\item Kenelm Digby Best, C.O., \textit{A Priest's Poems} (Burns \& Oates; Benziger Brothers, 1902), p.~136, for the public-domain English short \textit{Flos Carmeli}. Best supplies only the first two stanzas; the separate complete-English search records the historical near misses and three located modern complete versions.
\item Catholic Church, \textit{Proper of the Liturgy of the Hours of the Order of the Brothers of the Blessed Virgin Mary of Mount Carmel and the Order of Discalced Carmelites}, rev. and augm. ed. (Institutum Carmelitanum, 1993), \sourceurl{https://search.worldcat.org/title/41304501}{catalog record}. The complete Latin and English are reproduced from a \sourceurl{https://carmelitequotes.blog/2024/07/06/olmcnovena24-int/}{Carmelite-source transcription} that attributes the English to this Proper; Sullivan's article independently prints the first six English units. The transcription uses the plural \textit{nos ad esse / tecum in saeculum} reading and ends without \textit{Amen}. Its translator is not separately named, and the 1993 printed book was not directly inspected. The requester confirmed approval to use this and the two other located complete versions; the 1993 witness was selected for its joint institutional provenance and coherent Latin--English pairing. Sister Mary Paula Hayes, OCD, comp., \textit{Carmen Laudis: Hymns for Carmelite Feasts} (1984), represented by a separate \sourceurl{https://sacredheartretreathouse.com/blog/flos-carmeli/}{Carmelite-source transcription}, and the familiar 2007 version remain documented alternatives. The complete search and selection record is maintained with the repository research.
\end{itemize}

\end{multicols}
